\section{三种动力学绘景}
\subsection{连续表象中的算符核}
在进入三种绘景之前，先整理照片中的连续表象例题。算符 $\hat F$ 在坐标表象中的核为
\begin{equation}
F\left(x,x'\right)=\Braket{x|\hat F|x'},
\qquad
\left(\hat F\psi\right)\left(x\right)
=\int F\left(x,x'\right)\psi\left(x'\right)\,\mathrm{d}x'.
\end{equation}
必须保留算符次序，并明确导数作用于哪个变量。例如
\begin{equation}
\Braket{x|\hat x\hat p\hat x|x'}
=-\mathrm{i}\hbar xx'\frac{\partial}{\partial x}\delta\left(x-x'\right).
\end{equation}
这里左侧的 $\hat x$ 给出 $x$，右侧的 $\hat x$ 给出 $x'$；不能在导数尚未说明作用范围时把二者都直接写成同一个 $x$。

再以 $\hat F=\hat x^2\hat p$ 为例。在坐标表象与动量表象中分别有
\begin{align}
\Braket{x|\hat F|x'}
&=-\mathrm{i}\hbar x^2\frac{\partial}{\partial x}\delta\left(x-x'\right),\\
\Braket{p|\hat F|p'}
&=-\hbar^2p'\frac{\partial^2}{\partial p^2}\delta\left(p-p'\right).
\end{align}
第二式使用 $\Braket{p|\hat x|p'}=\mathrm{i}\hbar\partial_p\delta\left(p-p'\right)$，而最右侧的 $\hat p$ 作用在 $\Ket{p'}$ 上产生本征值 $p'$。这些核应作为分布理解，只有放入对测试波函数的积分后，分部积分等操作才具有明确意义。

\subsection{Schr\"odinger 绘景}
Schr\"odinger 绘景中，体系物理态按 Schr\"odinger 方程随时间演化，而没有显含时间的算符保持不变：
\begin{equation}
\Ket{\psi_{\mathrm{S}}\left(t\right)}=\hat U\left(t,t_0\right)
\Ket{\psi_{\mathrm{S}}\left(t_0\right)},
\qquad \hat A_{\mathrm{S}}=\hat A.
\end{equation}
对可能显含时间的 $\hat A_{\mathrm{S}}\left(t\right)$，期望值的导数为
\begin{equation}
\frac{\mathrm{d}}{\mathrm{d}t}\Braket{\hat A}_{\mathrm{S}}
=\frac{1}{\mathrm{i}\hbar}\Braket{\left[\hat A_{\mathrm{S}},\hat H_{\mathrm{S}}\right]}_{\mathrm{S}}
+\Braket{\frac{\partial\hat A_{\mathrm{S}}}{\partial t}}_{\mathrm{S}}.
\label{eq:expectation-evolution}
\end{equation}
因此，若算符不显含时间且与 Hamilton 算符对易，它所代表的物理量守恒。

\subsection{Heisenberg 绘景}
先定义统一的绘景变换，再区分体系物理态与算符本征态。取 $t_0$ 为绘景参考时刻，令 $\hat U\left(t,t_0\right)$ 为体系 Hamilton 算符生成的演化算符，且 $\hat U\left(t_0,t_0\right)=\hat I$。任意态矢量及算符在两种绘景之间按如下规则对应：
\begin{align}
\Ket{\phi_{\mathrm{H}}\left(t\right)}
&=\hat U^\dag\left(t,t_0\right)\Ket{\phi_{\mathrm{S}}\left(t\right)},\\
\hat A_{\mathrm{H}}\left(t\right)
&=\hat U^\dag\left(t,t_0\right)\hat A_{\mathrm{S}}\left(t\right)
\hat U\left(t,t_0\right).
\end{align}
第一式是态矢量的绘景变换，不要求任意给定的态矢量族都满足体系的 Schr\"odinger 方程。因此，“态不随时间变化”不是这个定义本身的结论，而是将它用于体系实际演化的物理态后得到的动力学结果。具体地，若
\begin{equation}
\Ket{\psi_{\mathrm{S}}\left(t\right)}
=\hat U\left(t,t_0\right)\Ket{\psi_{\mathrm{S}}\left(t_0\right)},
\end{equation}
则
\begin{equation}
\Ket{\psi_{\mathrm{H}}\left(t\right)}
=\hat U^\dag\left(t,t_0\right)\hat U\left(t,t_0\right)
\Ket{\psi_{\mathrm{S}}\left(t_0\right)}
=\Ket{\psi_{\mathrm{S}}\left(t_0\right)}
\equiv\Ket{\psi_{\mathrm{H}}}.
\end{equation}
这说明体系物理态在海森堡绘景中固定，而不意味着该绘景中的每一个 ket 都与时间无关。

对算符本征态也使用同一个变换。设不显含时间的 $\hat A$ 有一组离散正交完备本征态，满足 $\hat A\Ket{a_i}=a_i\Ket{a_i}$，选定其相位后定义
\begin{equation}
\Ket{a_i,t}_{\mathrm{H}}
=\hat U^\dag\left(t,t_0\right)\Ket{a_i}.
\end{equation}
由此可验证
\begin{equation}
\hat A_{\mathrm{H}}\left(t\right)\Ket{a_i,t}_{\mathrm{H}}
=\hat U^\dag\left(t,t_0\right)\hat A\Ket{a_i}
=a_i\Ket{a_i,t}_{\mathrm{H}}.
\end{equation}
幺正变换保持正交性与完备性，故固定的体系物理态可在这组随时间变化的基中展开：
\begin{equation}
\Ket{\psi_{\mathrm{H}}}
=\sum_i\Ket{a_i,t}_{\mathrm{H}}
\,{}_{\mathrm{H}}\Braket{a_i,t|\psi_{\mathrm{H}}}.
\end{equation}
基矢和展开系数一般都随时间变化，它们的变化共同保证左侧态矢量不变。这里的本征态族并非同一个初态按体系 Hamilton 算符正向演化所得的物理态族，因此不适用上面的时间不变结论。这与课堂照片中先定义海森堡态矢量、再讨论 $A$ 表象的顺序一致。

连续位置基按同样方式处理，所有等式按广义本征态的意义理解：
\begin{equation}
\Ket{x,t}_{\mathrm{H}}=\hat U^\dag\left(t,t_0\right)\Ket{x},
\qquad
\hat x_{\mathrm{H}}\left(t\right)\Ket{x,t}_{\mathrm{H}}
=x\Ket{x,t}_{\mathrm{H}}.
\end{equation}
在参考时刻，由于 $\hat U\left(t_0,t_0\right)=\hat I$，有
\begin{equation}
\Ket{x_0,t_0}_{\mathrm{H}}=\Ket{x_0}.
\end{equation}
这个等式来自共同的参考时刻和上述相位约定；仅凭本征方程不能固定本征态的相位。位置波函数是固定物理态在变化的位置基中的分量：
\begin{equation}
\psi\left(x,t\right)
={}_{\mathrm{H}}\Braket{x,t|\psi_{\mathrm{H}}}
=\Braket{x|\hat U\left(t,t_0\right)|\psi_{\mathrm{S}}\left(t_0\right)}.
\end{equation}
两种绘景给出相同的物理预言：
\begin{equation}
\Braket{\psi_{\mathrm{S}}\left(t\right)|\hat A_{\mathrm{S}}\left(t\right)|
\psi_{\mathrm{S}}\left(t\right)}
=\Braket{\psi_{\mathrm{H}}|\hat A_{\mathrm{H}}\left(t\right)|\psi_{\mathrm{H}}}.
\end{equation}
对定义式求导，得到 Heisenberg 运动方程
\begin{equation}
\boxed{\frac{\mathrm{d}\hat A_{\mathrm{H}}}{\mathrm{d}t}
=\frac{1}{\mathrm{i}\hbar}\left[\hat A_{\mathrm{H}},\hat H_{\mathrm{H}}\right]
+\left(\frac{\partial\hat A}{\partial t}\right)_{\mathrm{H}}.}
\label{eq:heisenberg-equation}
\end{equation}
其中
\begin{equation}
\hat H_{\mathrm{H}}\left(t\right)=\hat U^\dagger\left(t,t_0\right)
\hat H_{\mathrm{S}}\left(t\right)\hat U\left(t,t_0\right).
\end{equation}
若 $\hat H$ 不含时，则它与自身生成的 $\hat U$ 对易，故 $\hat H_{\mathrm{H}}=\hat H$。

\begin{example}[一维谐振子]
对
\begin{equation}
\hat H=\frac{\hat p^2}{2m}+\frac12m\omega^2\hat x^2,
\end{equation}
式\eqref{eq:heisenberg-equation}给出
\begin{equation}
\dot{\hat x}_{\mathrm{H}}=\frac{\hat p_{\mathrm{H}}}{m},
\qquad
\dot{\hat p}_{\mathrm{H}}=-m\omega^2\hat x_{\mathrm{H}}.
\end{equation}
所以算符本身满足经典形式的振动方程，解为
\begin{align}
\hat x_{\mathrm{H}}\left(t\right)
&=\hat x\cos\left[\omega\left(t-t_0\right)\right]
+\frac{\hat p}{m\omega}\sin\left[\omega\left(t-t_0\right)\right],\\
\hat p_{\mathrm{H}}\left(t\right)
&=\hat p\cos\left[\omega\left(t-t_0\right)\right]
-m\omega\hat x\sin\left[\omega\left(t-t_0\right)\right].
\end{align}
等价地，利用 $\left[\hat H,\hat a\right]=-\hbar\omega\hat a$，
\begin{equation}
\hat a_{\mathrm{H}}\left(t\right)=
\mathrm{e}^{-\mathrm{i}\omega\left(t-t_0\right)}\hat a,
\qquad
\hat a_{\mathrm{H}}^\dagger\left(t\right)=
\mathrm{e}^{\mathrm{i}\omega\left(t-t_0\right)}\hat a^\dagger.
\end{equation}
\end{example}

\subsection{相互作用绘景}
设 Hamilton 算符分解为
\begin{equation}
\hat H\left(t\right)=\hat H_0+\hat V\left(t\right),
\end{equation}
其中 $\hat H_0$ 不含时且其演化可精确求解。令
\begin{equation}
\hat U_0\left(t,t_0\right)=
\mathrm{e}^{-\mathrm{i}\hat H_0\left(t-t_0\right)/\hbar},
\end{equation}
并定义
\begin{align}
\Ket{\psi_{\mathrm{I}}\left(t\right)}
&=\hat U_0^\dagger\left(t,t_0\right)\Ket{\psi_{\mathrm{S}}\left(t\right)},\\
\hat A_{\mathrm{I}}\left(t\right)
&=\hat U_0^\dagger\left(t,t_0\right)\hat A_{\mathrm{S}}\left(t\right)
\hat U_0\left(t,t_0\right),\\
\hat V_{\mathrm{I}}\left(t\right)
&=\hat U_0^\dagger\left(t,t_0\right)\hat V\left(t\right)
\hat U_0\left(t,t_0\right).
\end{align}
直接求导可得
\begin{align}
\mathrm{i}\hbar\frac{\mathrm{d}}{\mathrm{d}t}\Ket{\psi_{\mathrm{I}}\left(t\right)}
&=\hat V_{\mathrm{I}}\left(t\right)\Ket{\psi_{\mathrm{I}}\left(t\right)},
\label{eq:interaction-state}\\
\frac{\mathrm{d}\hat A_{\mathrm{I}}}{\mathrm{d}t}
&=\frac{1}{\mathrm{i}\hbar}\left[\hat A_{\mathrm{I}},\hat H_0\right]
+\left(\frac{\partial\hat A}{\partial t}\right)_{\mathrm{I}}.
\label{eq:interaction-operator}
\end{align}
因此态由相互作用项 $\hat V_{\mathrm{I}}$ 驱动，算符的隐含时间演化由自由部分 $\hat H_0$ 驱动。期望值仍不变：
\begin{equation}
\Braket{\psi_{\mathrm{I}}\left(t\right)|\hat A_{\mathrm{I}}\left(t\right)|
\psi_{\mathrm{I}}\left(t\right)}
=\Braket{\psi_{\mathrm{S}}\left(t\right)|\hat A_{\mathrm{S}}\left(t\right)|
\psi_{\mathrm{S}}\left(t\right)}.
\end{equation}

定义相互作用绘景演化算符
\begin{equation}
\Ket{\psi_{\mathrm{I}}\left(t\right)}=\hat U_{\mathrm{I}}\left(t,t_0\right)
\Ket{\psi_{\mathrm{I}}\left(t_0\right)}.
\end{equation}
由式\eqref{eq:interaction-state}，
\begin{align}
\hat U_{\mathrm{I}}\left(t,t_0\right)
&=\mathcal T\exp\left[-\frac{\mathrm{i}}{\hbar}
\int_{t_0}^{t}\hat V_{\mathrm{I}}\left(t'\right)\,\mathrm{d}t'\right],
\label{eq:interaction-U}\\
\hat U_{\mathrm{I}}\left(t,t_0\right)
&=\hat U_0^\dagger\left(t,t_0\right)\hat U\left(t,t_0\right).
\end{align}
更一般地，
\begin{equation}
\hat U_{\mathrm{I}}\left(t,t'\right)
=\hat U_0^\dagger\left(t,t_0\right)\hat U\left(t,t'\right)
\hat U_0\left(t',t_0\right).
\end{equation}

\subsection{一阶微扰与线性响应}
把式\eqref{eq:interaction-U}保留到 $\hat V$ 的一阶：
\begin{equation}
\hat U_{\mathrm{I}}\left(t,t_0\right)
=\hat I-\frac{\mathrm{i}}{\hbar}\int_{t_0}^{t}
\hat V_{\mathrm{I}}\left(t'\right)\,\mathrm{d}t'+O\left(\hat V^2\right).
\end{equation}
若初态密度算符为 $\hat\rho_0$，则观测量的期望值到一阶为
\begin{align}
\Braket{\hat A\left(t\right)}
&=\Tr\left[\hat\rho_0\hat A_{\mathrm{I}}\left(t\right)\right]
+\frac{\mathrm{i}}{\hbar}\int_{t_0}^{t}\mathrm{d}t'
\Tr\left\{\hat\rho_0
\left[\hat V_{\mathrm{I}}\left(t'\right),\hat A_{\mathrm{I}}\left(t\right)\right]\right\}
+O\left(\hat V^2\right).
\end{align}
这一对易子形式是一阶线性响应理论的出发点。
